\chapter{PENGUJIAN DAN EVALUASI PRODUK}

\section{Metode Pengujian}

Pengujian dilakukan dengan beberapa pendekatan:

\begin{enumerate}
    \item \textbf{Functional Testing:} Memastikan aplikasi dapat diakses dan berfungsi dengan baik.
    \item \textbf{Security Testing:} Memeriksa konfigurasi keamanan jaringan dan kontainer.
    \item \textbf{Deployment Testing:} Menguji proses deployment dan rollback otomatis.
\end{enumerate}

\section{Skenario Pengujian dan Hasil}

\subsection{Pengujian Deployment}

Pengujian deployment dilakukan dengan membuat perubahan pada repositori Git dan mengamati sinkronisasi otomatis oleh ArgoCD.

\subsection{Pengujian Rollback}

Fitur rollback diuji dengan menggunakan histori revisi ArgoCD.

\section{Evaluasi Keamanan}

Aspek keamanan dievaluasi berdasarkan beberapa indikator:

\begin{table}[H]
\centering
\caption{Evaluasi Keamanan}
\begin{tabular}{|l|l|l|}
\hline
\textbf{Aspek} & \textbf{Implementasi} & \textbf{Status} \\
\hline
Akses Eksternal & Cloudflare Tunnel (outbound-only) & \checkmark \\
Network Isolation & Default-deny NetworkPolicies & \checkmark \\
Secret Management & Sealed Secrets (RSA encrypted) & \checkmark \\
Pod Security & Non-root, read-only filesystem & \checkmark \\
TLS & cert-manager (auto-renewal) & \checkmark \\
\hline
\end{tabular}
\end{table}

\section{Analisis Hasil}

Berdasarkan hasil implementasi dan pengujian, sistem yang dibangun telah memenuhi tujuan penelitian.

\section{Kesimpulan dan Saran}

\subsection{Kesimpulan}

Berdasarkan penelitian yang telah dilakukan, dapat diambil kesimpulan sebagai berikut:

\begin{enumerate}
    \item Arsitektur GitOps untuk deployment aplikasi pada klaster Kubernetes berhasil dirancang dengan komponen utama meliputi repositori Git, ArgoCD, dan komponen infrastruktur pendukung.
    \item Implementasi GitOps menggunakan ArgoCD untuk deployment aplikasi InvenioRDM berhasil dilakukan dengan mekanisme sinkronisasi otomatis dan pengaturan urutan deployment melalui sync waves.
    \item Sistem yang dibangun menunjukkan efektivitas dalam deployment otomatis dan kemudahan rollback. Aspek keamanan juga terpenuhi.
\end{enumerate}

\subsection{Saran}

Berdasarkan hasil penelitian, saran untuk pengembangan selanjutnya adalah:

\begin{enumerate}
    \item Menambahkan mekanisme canary deployment atau blue-green deployment.
    \item Mengevaluasi penggunaan GitOps controller alternatif seperti Flux CD.
    \item Mengimplementasikan solusi secret management yang lebih advanced seperti HashiCorp Vault.
    \item Melakukan pengujian load testing dan chaos engineering.
\end{enumerate}
